\section{Discussion}

\subsection{Interpretation of Validation Results}

Phase accuracy to $0.2^\circ$ corresponds to 4 picosecond delay errors at L-band, well below atmospheric turbulence effects (nanoseconds). Amplitude measurements within 1\% confirm proper signal integration. Closure phases near zero validate the entire processing chain independently of calibration. Noise scaling as $\sigma_V \propto 1/\sqrt{\Delta t \Delta\nu}$ confirms proper coherent signal accumulation. Multi-source validation demonstrates correct superposition through electromagnetic field linearity.

\subsection{Limitations and Scientific Applications}

Key limitations: narrow field of view (single phase centre), fixed spectral resolution, no calibration pipeline, limited polarization (XX/YY only), marginal real-time performance for large arrays, and data format compatibility requiring conversion for CASA/AIPS.

Despite limitations, the correlator suits continuum source observation, broad spectral line studies, time-domain astronomy (1.7$\times$ real-time enables transient searches), and interferometry education. For Nooitgedacht's L-band demonstration goals, the system provides appropriate functionality.

\subsection{Comparison and Validation Confidence}

Compared to established correlators (DiFX for VLBI, xGPU for large arrays, HPGBE for moderate scales), the Python implementation sacrifices raw performance for simplicity and transparency, appropriate for teaching and small-array prototyping.

Validation caveats: testing used simulated data (real telescope data contains RFI and artifacts), limited parameter space (2–8 antennas, 64–1024 channels), and lacks cross-validation with DiFX. While validation demonstrates correctness for tested scenarios, on-sky observations would strengthen confidence.

Key lessons: test-driven development accelerated debugging; analytical validation provided unambiguous criteria; profiling revealed FFT (not correlation) dominated; double precision proved necessary for phase accuracy; multi-tier testing caught errors at appropriate levels.

Future testing directions include on-sky observations, long-duration stability tests, comparative benchmarking with DiFX, stress testing with corrupted inputs, and expanding regression suites for new features. The testing establishes the correlator as functional and accurate for Nooitgedacht's research objectives.
